\Article{23}
Entwickelt man die Potenz $(x^4+1)^{\frac12(p-1)}$ nach dem Binomialtheorem, so treten drei Glieder auf, in denen der Exponent von $x$ durch $p-1$ teilbar ist, nämlich
\[
x^{2(p-1)},\quad Px^{p-1}\quad\text{und }1
\]
\SourceLeaf{100}%
wobei $P$ den mittleren Koeffizienten
\[
\frac{\frac12(p-1)\pdot\frac12(p-3)\pdot\frac12(p-5)\dotsfive\frac14(p+3)}{1\pdot2\pdot3\dotsfive\frac14(p-1)}
\]
bezeichnet. Setzen wir also für $x$ der Reihe nach die Zahlen $1,2,3\dotsfour p-1$ ein, so erhalten wir nach dem Hilfssatz in Artikel 19
\[
\Sigma(x^4+1)^{\frac12(p-1)}\equiv-2-P
\]
Beachten wir aber das in Artikel 19 Dargelegte und außerdem, dass die Zahlen der Sammlungen $A,B,C,D$, zur Potenz mit dem Exponenten $\frac12(p-1)$ erhoben, modulo $p$ jeweils zu den Zahlen $+1,-1,+1,-1$ kongruent sind, so sehen wir leicht, dass
\[
\Sigma(x^4+1)^{\frac12(p-1)}\equiv4(00)-4(01)+4(02)-4(03)
\]
und somit nach den Schemata am Ende der Artikel 18, 20
\[
\Sigma(x^4+1)^{\frac12(p-1)}\equiv-2a-2
\]
gilt. Der Vergleich dieser beiden Werte liefert einen besonders eleganten Satz: Wir haben
\[
P\equiv2a\ (\mathrm{mod.}\ p)
\]

Bezeichnen wir die vier Produkte
\[
\begin{array}{l}
1\pdot2\pdot3\mathpunct{.\,.\,.\,.\,.\,.\,.\,.\,.\,.\,.\,.\,.\,.\,.\,.\,.\,.\,.\,.\,.\,.\,.\,.\,.}\frac14(p-1)\\
\frac14(p+3)\pdot\frac14(p+7)\pdot\frac14(p+11)\dotsfive\frac12(p-1)\\
\frac12(p+1)\pdot\frac12(p+3)\pdot\frac12(p+5)\dotsfive\frac34(p-1)\\
\frac14(3p+1)\pdot\frac14(3p+5)\pdot\frac14(3p+9)\dotsfive(p-1)
\end{array}
\]
jeweils mit $q,r,s,t$, so lässt sich der vorangegangene Satz als
\[
2a\equiv\frac rq\ (\mathrm{mod.}\ p)
\]
schreiben.

\Needspace{18\baselineskip}
Da jeder Faktor von $q$ seine Ergänzung zu $p$ in $t$ hat, gilt $q\equiv t\ (\mathrm{mod.}\ p)$, wenn die Anzahl der Faktoren gerade ist, also wenn $p$ die Form $8m+1$ hat, dagegen $q\equiv-t$, wenn die Anzahl der Faktoren ungerade ist oder $p$ die Form $8m+5$ hat. Ebenso gilt im ersten Fall $r\equiv s$, im zweiten $r\equiv-s$. In beiden Fällen gilt $qr\equiv st$, und da bekanntlich $qrst\equiv-1$, gilt $qqrr\equiv-1$,
\SourceLeaf{101}%
und somit $qr\equiv\pm f\ (\mathrm{mod.}\ p)$. Verbinden wir diese Kongruenz mit dem eben gefundenen Satz, so erhalten wir $rr\equiv\pm2af$ und daher nach den Artikeln \mbox{19, 20}
\[
2b\equiv\pm rr\ (\mathrm{mod.}\ p)\text{*)}
\]
\begingroup\renewcommand{\thefootnote}{*)}\footnotetext{und $\{(a\mp b)q\}^2\equiv a\equiv\left(\frac{r-qrr}{2}\right)^2$}\endgroup

Es ist sehr bemerkenswert, dass sich die Zerlegung der Zahl $p$ in zwei Quadrate durch völlig direkte Rechenoperationen finden lässt: Die Wurzel des ungeraden Quadrats ist der absolut kleinste Rest von $\frac r{2q}$, die Wurzel des geraden Quadrats dagegen der absolut kleinste Rest von $\frac12rr$ zum Modul $p$. Den Ausdruck $\frac r{2q}$, der für $p=5$ den Wert $=1$ hat, kann man für größere Werte von $p$ auch so schreiben:
\[
\frac{6\pdot10\pdot14\pdot18\dotsfive(p-3)}{2\pdot3\pdot4\pdot5\dotsfive\frac14(p-1)}
\]

Da wir aber außerdem wissen, mit welchem Vorzeichen diese Formel die Wurzel des ungeraden Quadrats liefert, nämlich mit demjenigen, durch das sie stets die Form $4m+1$ hat, verdient es besondere Aufmerksamkeit, dass bisher kein ähnliches allgemeines Kriterium für das Vorzeichen der Wurzel des geraden Quadrats gefunden werden konnte. Wer ein solches findet und uns mitteilt, dem werden wir sehr dankbar sein. Einstweilen erschien es angebracht, hier die Werte der Zahlen $a,b,f$ anzufügen, die sich für Werte von $p$ unter $200$ aus den kleinsten Resten der Ausdrücke $\frac r{2q}$, $\frac12rr$, $qr$ ergeben.

\SourceLeaf{102}%
\[
\begin{array}{r|r|r|r}
p&a&b&f\\\hline
5&+1&+2&2\\
13&-3&-2&5\\
17&+1&-4&13\\
29&+5&+2&12\\
37&+1&-6&31\\
41&+5&+4&9\\
53&-7&-2&23\\
61&+5&-6&11\\
73&-3&-8&27\\
89&+5&-8&34\\
97&+9&+4&22\\
101&+1&-10&91\\
109&-3&+10&33\\
113&-7&+8&15\\
137&-11&+4&37\\
149&-7&-10&44\\
157&-11&-6&129\\
173&+13&+2&80\\
181&+9&+10&162\\
193&-7&+12&81\\
197&+1&-14&183
\end{array}
\]
\BookRule{.28\linewidth}
